% pyALDIC Quick Guide -- two-page, two-column reference sheet.
% Compile the same way as user-guide.tex (pdflatex / xelatex / latexmk).

\documentclass[10pt,letterpaper,twocolumn]{article}

\usepackage[T1]{fontenc}
\usepackage[utf8]{inputenc}
\usepackage{lmodern}

\usepackage[margin=0.55in,columnsep=0.3in,top=0.5in,bottom=0.5in]{geometry}
\usepackage{xcolor}
\definecolor{aldicblue}{RGB}{65,105,168}
\definecolor{aldicgreen}{RGB}{34,150,94}
\definecolor{aldicred}{RGB}{192,69,72}

\usepackage[colorlinks=true,linkcolor=aldicblue,urlcolor=aldicblue]{hyperref}

\usepackage{enumitem}
\setlist{topsep=1pt,itemsep=0pt,parsep=0pt,leftmargin=*}

\usepackage{titlesec}
\titlespacing*{\section}{0pt}{4pt}{1pt}
\titlespacing*{\subsection}{0pt}{3pt}{0pt}
\titleformat{\section}{\normalsize\bfseries\color{aldicblue}}{\thesection.}{4pt}{}
\titleformat{\subsection}{\small\bfseries}{}{0pt}{}

\usepackage{booktabs}
\usepackage{array}

\setlength{\parindent}{0pt}
\setlength{\parskip}{2pt}
\pagestyle{empty}

\begin{document}

\twocolumn[%
\begin{center}
  {\Large\bfseries\color{aldicblue} pyALDIC Quick Guide}\\[2pt]
  {\small Two-page reference. Full walkthrough: \texttt{user-guide.tex}.}
\end{center}
\vspace{6pt}
]

\section{Launch}

\texttt{al-dic} (or \texttt{python -m al\_dic}).

\section{Load images}

Drop an image folder onto the drop zone or click Browse. All files
must be the same size. Frame 1 is always the reference.

\section{Choose Workflow Type}

This decides which frames need a Region of Interest, so do it
\emph{before} drawing.

{\small
\begin{tabular}{@{}p{0.42\linewidth}p{0.42\linewidth}@{}}
  \toprule
  Experiment & Tracking + Ref Update \\
  \midrule
  Soft material $<$2\% strain & Accumulative \\
  2--30\% strain, any rotation $>$5$^\circ$, impact & Incremental + Every Frame \\
  Known quasi-static load steps & Incremental + Custom Frames \\
  \bottomrule
\end{tabular}
}

Solver: \textbf{AL-DIC} default (smoother, $\sim 3\times$ slower);
\textbf{Local DIC} for quick iteration.

\section{Draw Region of Interest}

The blue-accent hint above the toolbar tells you which frames need
a region. Toolbar:

\begin{itemize}
  \item \textbf{+ Add}, \textbf{Cut}, \textbf{+ Refine} (brush,
        frame 1 only). Each opens a Polygon / Rectangle / Circle
        menu.
  \item \textbf{Import} / \textbf{Batch Import} -- load PNG masks.
  \item \textbf{Save} / \textbf{Invert} / \textbf{Clear} -- per-
        frame mask operations.
\end{itemize}

For per-frame editing, click the \textcolor{aldicred}{Need} /
\textcolor{aldicgreen}{Edit} / Add button on a row in the image
list.

\section{Parameters}

{\small
\begin{tabular}{@{}p{0.28\linewidth}p{0.58\linewidth}@{}}
  \toprule
  Control & Guidance \\
  \midrule
  Subset Size & 41 default; larger (51--81) for noise, smaller
                (21--31) for fine speckle. \\
  Subset Step & 8 for 512\textsuperscript{2}, 16 for
                1024\textsuperscript{2}. Halving step $\to$
                $\sim 4\times$ runtime. \\
  Search Range & 20 default; raise for large inter-frame motion.
                 Auto-expand is on. \\
  Refinement & Check Inner / Outer to densify mesh near hole or
               ROI boundaries; pick a level (Light to Ultra). \\
  \bottomrule
\end{tabular}
}

\textbf{Advanced} section (collapsed by default) holds the initial-
guess mode and auto-expand toggle; defaults are fine.

\section{Run \& Cancel}

Click the blue \textbf{Run DIC Analysis}. Progress, elapsed /
remaining, and console log appear on the right. Red
\textbf{Cancel} stops cleanly (partial results kept).

Fatal errors raise a modal dialog; the full traceback is in the
console log.

\section{View results}

On the right sidebar after a run:
\begin{itemize}
  \item \textbf{FIELD}: Disp U / Disp V toggle. \emph{Show on
        deformed frame} checkbox sits here (controls WHERE, not
        how).
  \item \textbf{VISUALIZATION}: Colormap, auto / manual range,
        opacity.
  \item \textbf{PHYSICAL UNITS}: Pixel size + frame rate -- scales
        displacement / velocity output; strain is dimensionless.
\end{itemize}

\section{Strain post-processing}

A separate \textbf{Strain window} opens automatically when the run
finishes.

\begin{itemize}
  \item \textbf{Method}: \emph{Plane fitting} (smooth, default) vs
        \emph{FEM nodal} (higher resolution, noisier).
  \item \textbf{VSG size}: Virtual strain gauge diameter
        (default 41 px).
  \item \textbf{Smoothing}: Off / Light (0.5$\times$ step) /
        Medium (1$\times$ step) / Strong (2$\times$ step)
        $\blacktriangle$.
  \item \textbf{Strain type}: \textbf{Infinitesimal} (small
        deformation only), \emph{Eulerian-Almansi},
        \textbf{Green-Lagrangian} (\textbf{use for any rotation
        $>$ 2$^\circ$}; it is exactly zero under rigid rotation).
\end{itemize}

Click green \textbf{Compute Strain} whenever parameters change.
Field selector now exposes \texttt{strain\_exx},
\texttt{strain\_eyy}, \texttt{strain\_exy}, principal strains,
\texttt{strain\_maxshear}, \texttt{strain\_von\_mises},
\texttt{strain\_rotation}, plus the displacement fields.

\section{Export}

Both windows have an \textbf{Export Results} button; they open the
same 4-tab dialog:

\begin{itemize}
  \item \textbf{Data}: NPZ / MAT / CSV of any selected fields.
  \item \textbf{Images}: per-field PNG / PDF per frame. Colorbar
        ON by default.
  \item \textbf{Animations}: MP4 / GIF. Requires \texttt{ffmpeg}
        on PATH for MP4 (falls back to GIF otherwise).
  \item \textbf{Report}: multi-page PDF summary of the run.
\end{itemize}

\section{Save / resume sessions}

\textbf{File $\to$ Save Session} (\texttt{Ctrl+S}) writes a single
\texttt{.aldic.json} with every Region of Interest (base64 PNG),
all parameters, and physical units. Does NOT save pipeline results
(re-run from Session is cheap).

\textbf{File $\to$ Open Session} (\texttt{Ctrl+O}) restores
everything. If the image folder is gone, parameters and Regions of
Interest still load; re-drop the images manually.

\section{Troubleshooting quick table}

{\small
\begin{tabular}{@{}p{0.42\linewidth}p{0.50\linewidth}@{}}
  \toprule
  Symptom & Likely cause / fix \\
  \midrule
  Run button disabled &
    No images or no Region of Interest on frame 1. \\
  Crazy strain under big rotation &
    Using Infinitesimal strain; switch to Green-Lagrangian. \\
  FFT peaks clipped warnings &
    Raise Search Range or keep auto-expand on. \\
  Refine brush grayed out &
    Brush is frame-1 only; navigate to frame 1. \\
  Results poor near hole edges &
    Enable \emph{Refine Inner Boundary}. \\
  Slow run &
    Halve Subset Step for next iteration;
    use Local DIC for previews. \\
  Session won't load &
    Check console: JSON corrupt, future schema,
    or missing image folder (warning only -- params restore). \\
  \bottomrule
\end{tabular}
}

\end{document}
